% Self-contained copy for the AAG tables file. tabularx fills the page width (no
% adjustbox: tabularx must not be wrapped in \resizebox). 11 rows fit the page height.
% Source of record: reports/tables/segmentation_variables.tex
\begin{table}[htbp]
\centering
\small
\caption{Grouped built environment categories and their constituent variables from street view imagery analysis.}
\label{tab:segmentation_variables}
\begin{tabularx}{\textwidth}{llX}
\toprule
\textbf{Category} & \textbf{Type} & \textbf{Constituent Variables} \\
\midrule
\multicolumn{3}{l}{\textit{Grouped categories (composite scores)}} \\
Greenery & Ratio & Vegetation \\
Open Space & Ratio & Sky, Terrain, Sand \\
Building & Ratio & Building \\
Road & Ratio & Road, Service Lane \\
Active Mobility Infrastructure & Ratio & Sidewalk, Bike Lane, Pedestrian Area \\
Active Mobility Presence & Count & Person, Bicycle, Bicyclist \\
Vehicle Presence & Count & Car, Truck, Bus, Motorcycle, Trailer, Other Vehicle, Caravan \\
Physical Boundaries & Ratio & Fence, Wall, Barrier \\
\midrule
\multicolumn{3}{l}{\textit{Individual features (not grouped)}} \\
Symbolic US Flag & Count & American Flag \\
CCTV Surveillance & Count & CCTV Camera \\
Visual Complexity & --- & Shannon entropy of all segmentation pixel ratios \\
\bottomrule
\end{tabularx}
\begin{flushleft}
\footnotesize
Ratio variables are derived from semantic segmentation (pixel area proportions); count variables from object detection. For each grouped category, constituent variables were z-scored and summed to create composite scores, then re-standardized. Individual features were z-scored directly.
\end{flushleft}
\end{table}
